% ===================================================================
%  TAULA N° 12 — L'Estacion. — Los Camins de fèrre. — Los Naviris.
%  Traduction 2026 d'apres texto/io, controlee sur texto/fr et
%  texto/ca. Consigne : voir texto/oc/10-taula-01.tex.
% ===================================================================

%%K t12-apar-1 apar
\VUsaut{4.0mm}
\VUtitre{15.0pt}{60}{TAULA N° 12}
\VUsaut{2.4mm}
\VUpk{11.4pt}{40}{L'Estacion. --- Los Camins de fèrre. --- Los Naviris.}
\VUsaut{3.4mm}
\textsc{Nòta.} --- \textit{Los oraris emplegats dins cada país son diferents, per
aquò avèm pres coma exemple las oras del} \VUgras{quadre francés}.

%%K t12-01-1 p
1. --- Veni de far un viatge per Alemanha. Abans de l'entreprene aviái crompat una
\VUgras{guida d'oraris} \textsuperscript{(15)} per saber l'ora de partença dels
trens. Aprèp aver verificat que lo tren mai comòde partiá de París a nòu oras del
vèspre e arribava a Estrasborg a una ora e tretze, preparèri mon viatge e,
l'endeman, me faguèri menar a l'estacion.

%%K t12-02-1 p
2. --- Quand dintrèri dins la granda sala de partenças, darrièr un vièlh
\VUgras{curat} \textsuperscript{(2)} de vilatge que portava a la man sa
\VUgras{saca de viatge} \textsuperscript{(3)}, èran ja arribats fòrça
\VUgras{viatjaires} \textsuperscript{(1)}.

%%K t12-02-2 p
Coma voliái mandar un \VUgras{telegrama} \textsuperscript{(5)}, faguèri qu'un
\VUgras{contrarotlaire} \textsuperscript{(6)} m'indiquèsse lo \VUgras{burèu de
telegraf} \textsuperscript{(4)}.

%%K t12-03-1 p
3. --- Abans de passar a la \VUgras{sala d'espèra} \textsuperscript{(7)} anèri
prene un \VUgras{bilhet} \textsuperscript{(9)} d'anada e retorn.

%%K t12-04-1 p
4. --- Esperèri pas gaire al \VUgras{guichet} \textsuperscript{(8)}, que i aviá
pauc de mond. Un \VUgras{soldat} \textsuperscript{(10)} que partiá en permission me
cedèt amablament son torn, de biais qu'aguèri lo temps de demandar qualques
entresenhas al \VUgras{cap d'estacion} \textsuperscript{(13)}, que parlava amb lo
\VUgras{comissari} \textsuperscript{(12)} de susvelhança e amb lo
\VUgras{gendarma} \textsuperscript{(11)} de servici.

%%K t12-tit-1 sub
\VUpk{10.2pt}{40}{Enregistrar los bagatges.}
\VUsaut{3.0mm}

%%K t12-05-1 p
5. --- Aprèp aver crompat un jornal al \VUgras{quiòsc de libres}
\textsuperscript{(14)}, faguèri pesar mos \VUgras{bagatges}
\textsuperscript{(17)}, qu'un \VUgras{portaire} \textsuperscript{(16)} veniá de
portar sus una \VUgras{bèrla} \textsuperscript{(19)}; l'\VUgras{emplegat}
\textsuperscript{(22)} encargat de la \VUgras{bascula} \textsuperscript{(21)} me
diguèt que ma mala pesava trenta-cinc quilòs, çò que fasiá un \VUgras{excès} de
cinc quilòs, per lo qual me caliá pagar un suplement al \VUgras{guichet dels
bagatges} \textsuperscript{(23)}.

%%K t12-06-1 p
6. --- Coma èri en avança, aguèri l'escasença d'agachar lo movement dels
viatjaires dins l'estacion. Los uns daissavan o recampavan de pichons paquets a la
\VUgras{consigna} \textsuperscript{(25)}; los autres consultavan los
\VUgras{oraris} \textsuperscript{(26)}. Los viatjaires qu'arribavan se despachavan
cap a la \VUgras{sortida} \textsuperscript{(32)} amb lor \VUgras{valisa}
\textsuperscript{(24)} a la man. D'unes esperavan a la \VUgras{liurason dels
bagatges} \textsuperscript{(27)} e presentavan lor \VUgras{recepissat}
\textsuperscript{(29)} per retirar lors malas, \VUgras{caissas}
\textsuperscript{(28)} o \VUgras{panièrs} \textsuperscript{(30)}.

%%K t12-07-1 p
7. --- Los \VUgras{emplegats de l'octròi} \textsuperscript{(33)} regardavan amb
suènh los paquets per veire s'i aviá quicòm a declarar.

%%K t12-08-1 p
8. --- D'unes viatjaires se dirigissián cap al \VUgras{bufet}
\textsuperscript{(34)}, ont un \VUgras{garçon} \textsuperscript{(35)} se despachava
de los servir. Lor caliá se despachar, que lo \VUgras{tren} \textsuperscript{(36)},
qu'èra un exprès, s'arrestava pas gaire de temps a l'estacion.

%%K t12-09-1 p
9. --- Passèri puèi sul cai de partença e cerquèri un \VUgras{vagon}
\textsuperscript{(37)} dirècte per pas aver de cambiar pendent lo viatge. Pugèri
dins una veitura a corredor qu'a un \VUgras{compartiment} \textsuperscript{(38)}
pels fumaires e un autre reservat a las «dònas solas».

%%K t12-tit-2 sub
\VUpk{10.2pt}{40}{Partença de nuèch.}
\VUsaut{3.0mm}

%%K t12-10-1 p
10. --- Aprèp aver pujat lo \VUgras{marchapè} \textsuperscript{(40)}, barrat la
\VUgras{portièra} \textsuperscript{(39)}, remontat lo \VUgras{ridèu}
\textsuperscript{(41)} e pausat mos pichons paquets dins lo \VUgras{filat}
\textsuperscript{(43)}, m'assetèri sul \VUgras{banc} \textsuperscript{(42)}. En
veire lo \VUgras{lampista} \textsuperscript{(44)} alucar los lums, pensèri qu'aviái
una nuèch a passar dins lo vagon e sonèri l'emplegat encargat del loguièr dels
\VUgras{coissins} \textsuperscript{(45)} per ne demandar un.

%%K t12-11-1 p
11. --- Lo conductor del tren venguèt verificar los bilhets. Li demandèri perqué
partissiam pas, qu'èra l'ora. Me respondèt que lo \VUgras{cap de tren}
\textsuperscript{(46)} èra ocupat a far pausar de bicicletas dins lo
\VUgras{furgon} \textsuperscript{(47)}. En mai, s'èran pas encara cambiats totes
los \VUgras{escaufapès} \textsuperscript{(48)}.

%%K t12-12-1 p
12. --- Mas anàvem pas tardar a partir, que lo \VUgras{caufaire}
\textsuperscript{(50)}, sus son \VUgras{tènder} \textsuperscript{(49)}, prenia de
carbon per revifar lo fuòc de la \VUgras{locomotiva} \textsuperscript{(54)}, e
\VUgras{lo mecanician} \textsuperscript{(51)} esperava pas que lo senhal del
\VUgras{sotscap d'estacion} \textsuperscript{(52)}, qu'èra sul \VUgras{cai}
\textsuperscript{(53)}.

%%K t12-13-1 p
13. --- La \VUgras{caudièra} \textsuperscript{(56)} de la locomotiva bramava
sordament; la \VUgras{chiminèa} \textsuperscript{(55)} vomissiá de rajòls de fum
mesclat de \VUgras{vapor} \textsuperscript{(60)}. Sonèt un \VUgras{siblet}
\textsuperscript{(59)} agut e los \VUgras{pistons} \textsuperscript{(58)} se
metèron en movement, en fasent virar las \VUgras{ròdas} \textsuperscript{(57)} de
la maquina pesuca.

%%K t12-tit-3 sub
\VUsaut{3.0mm}
\VUpk{10.2pt}{40}{En rota!}
\VUsaut{3.0mm}

%%K t12-14-1 p
14. --- La velocitat del tren, que coriá suls \VUgras{rails}
\textsuperscript{(64)}, s'accelerèt; los \VUgras{cais} \textsuperscript{(65)}
s'acabèron; passèrem los \VUgras{comuns} \textsuperscript{(66)} e tanben una
seguida de vagons apartats sus una autra \VUgras{via} \textsuperscript{(63)}: de
\VUgras{veituras-lièch} \textsuperscript{(67)}, una \VUgras{veitura-restaurant}
\textsuperscript{(68)}, una \VUgras{veitura-pòsta} \textsuperscript{(69)}.

%%K t12-noto-1 noto
\VUnotes{143.2mm}{%
\textit{(*) Nòta.} --- \textit{S'entend que lo viatge es descrich tal coma èra al
long de la frontièra d'abans la guèrra.}%
}

%%K t12-15-1 p
15. --- Puèi passèt davant nòstres uèlhs l'\VUgras{estacion de mercadarias}
\textsuperscript{(71)}. Jos los cobèrts, d'emplegats cargavan las
\VUgras{mercadarias} \textsuperscript{(72)} mandadas a pichona velocitat. De
\VUgras{manòbras} \textsuperscript{(76)}, contra lo \VUgras{depaus
d'aiga} \textsuperscript{(73)}, manobravan de \VUgras{vagons de bestial}
\textsuperscript{(75)} e de \VUgras{vagons plataforma} \textsuperscript{(79)} per
formar un \VUgras{tren de mercadarias} \textsuperscript{(80)}.

%%K t12-16-1 p
16. --- Passèrem la darrièra \VUgras{agulha} \textsuperscript{(78)}, contra la
quala èra l'agulhaire; daissèrem darrièr lo \VUgras{disc de senhals}
\textsuperscript{(86)} e l'estacion se perdèt al luènh.

%%K t12-17-1 p
17. --- Lèu dintrèrem sul \VUgras{pont de fèrre} \textsuperscript{(74)} que crosa
lo \VUgras{riu} \textsuperscript{(81)}. Passat lo pont, seguiguèrem lo cors del
riu, ont de poderoses \VUgras{remolcadors} \textsuperscript{(82)} tiravan de
pesucas \VUgras{gabarras} \textsuperscript{(83)}. Vegèrem en mai un \VUgras{naviri
de passatgièrs} \textsuperscript{(84)} al moment ont anava acostar al
\VUgras{ponton} \textsuperscript{(85)}.

%%K t12-18-1 p
18. --- Aprèp aver passat de long mantuna estacion, traversèrem qualques
\VUgras{tunèls} \textsuperscript{(87)} e daissèrem darrièr nosautres
d'innombrablas \VUgras{cabanas} \textsuperscript{(88)} de
\VUgras{gardabarrièras}. Bressat pel secodiment del tren, m'endormiguèri
prigondament.

%%K t12-tit-4 sub
\VUsaut{3.0mm}
\VUpk{10.2pt}{40}{L'estacion de Deutsch-Avricourt.}
\VUsaut{3.0mm}

%%K t12-19-1 p
19. --- Me desrevelhèt tot d'un còp la votz d'un emplegat que cridava:
«Deutsch-Avricourt! \textsuperscript{(*)} Davalatz totes!». Èra l'inspeccion de
las doanas e totas las fastidiosas formalitats de la frontièra. Me calguèt pausar
mon relòtge a l'ora amb lo \VUgras{relòtge} \textsuperscript{(89)} de l'estacion,
qu'anava cinquanta-cinc minutas en avança; a pena desrevelhat, destriavi amb pena
l'\VUgras{agulha granda} \textsuperscript{(90)} de la \VUgras{pichona}
\textsuperscript{(91)}; lo \VUgras{quadran} \textsuperscript{(92)} el meteis èra
tot confús per la nèbla del matin.

%%K t12-20-1 p
20. --- Mentre que los doanièrs fasián lor inspeccion, deschifrèri en badalhar
las \VUgras{afichas} \textsuperscript{(93)} qu'ornan las parets de l'estacion.
Dejunèri per passar lo temps, consultèri la mapa de la \VUgras{ret}
\textsuperscript{(94)} alemanda per veire quina distància me separava encara
d'Estrasborg, e tornèri pujar dins lo vagon, reconfortat per l'esperança
d'arribar lèu al tèrme de mon viatge.
\VUsaut{6.0mm}
\VUfilet{22mm}
